\chapter{Module Federation}

The host shell loads \emph{everything} above the WASM iframe through
\texttt{@angular-architects/native-federation}. Adding a new
route-bearing remote means adding a row to \texttt{federation\_modules},
shipping a per-remote bundle, and pointing \texttt{loadComponent} at the
exposed module. Adding a new game means the same plus an
\texttt{<app-game>-remote>} Web Component the host can place in its
template.

\section{The two ordering rules}

\begin{enumerate}
  \item \textbf{\texttt{initFederation} before \texttt{bootstrapApplication}.}
        \snapfile{frontend/src/main.ts} dynamic-imports
        \snapfile{bootstrap.ts} after the federation runtime resolves.
        Static imports of \texttt{@angular/*} from \texttt{main.ts} would
        be evaluated before the import map is installed and would
        therefore fail to resolve.
  \item \textbf{Remote-loader access through \texttt{globalThis}.}
        Route handlers reach the federation runtime through
        \texttt{globalThis.\_\_SNAPDRAGON\_LOAD\_REMOTE\_MODULE\_\_}, set
        by \texttt{main.ts} once \texttt{initFederation} resolves. This
        keeps \texttt{@angular-architects/native-federation} out of
        karma's static graph and out of every route-handler bundle.
\end{enumerate}

\section{The manifest}

\begin{figure}[!h]
\centering
\snapdiagram{diagrams/federation-bootstrap.tex}
\caption{Federation bootstrap. The host fetches the manifest from the BFF
\emph{before} \texttt{bootstrapApplication}, then constructs the import
map, then mounts the host \texttt{AppComponent}. New remotes appear by
inserting rows into \texttt{federation\_modules} — no host redeploy
required.}
\end{figure}

The manifest endpoint is
\texttt{GET /api/federation/manifest} — DB-backed, public on purpose
because the host needs the \texttt{login} remote URL before any user has
authenticated. It returns the active rows of \texttt{federation\_modules}
ordered by \texttt{kind, slug}:

\begin{lstlisting}[language=json]
{
  "generatedAt": "2026-04-26T05:30:11Z",
  "remotes": [
    { "slug": "login",  "kind": "remote-shell",
      "remoteEntryUrl": "/_remotes/login/remoteEntry.json",
      "exposedModule": "./LoginComponent",
      "displayName": "Login",  "version": "1.0.0", "metadata": {} },
    { "slug": "admin",  "kind": "remote-shell",
      "remoteEntryUrl": "/_remotes/admin/remoteEntry.json",
      "exposedModule": "./AdminShell",
      "displayName": "Admin",  "version": "1.0.0", "metadata": {} },
    /* ...games-three-wizards, games-dodecahedron, wasm-fragment... */
  ]
}
\end{lstlisting}

\section{Shared dependencies}

\snapfile{frontend/federation-base.js} consolidates the shared-package
configuration so the host and all four remotes ship the same Angular
version:

\begin{lstlisting}[language=javascript]
// One source of truth for shareAll + skip helpers.
const STD_SHARED = ['@angular/core','@angular/common','@angular/forms',
                    '@angular/router','rxjs','zone.js'];
const STD_SKIP   = ['@angular/compiler','@angular-architects/native-federation'];

module.exports = {
  snapdragonHost:   () => snapdragonShared({singleton:true,strictVersion:true}),
  snapdragonRemote: () => snapdragonShared({singleton:true,requiredVersion:'auto'}),
};
\end{lstlisting}

The angular-architects builder enforces that \texttt{@angular/core} and
the rest of the standard set are loaded \emph{exactly once} per page;
mismatched versions across remotes produce the
\texttt{NG0203 inject() must be called from an injection context} runtime
error, which is the federation-shaped variant of the more familiar
\texttt{NullInjectorError}.

\section{Routing into a remote}

The \texttt{loadRemote<T>} helper in \snapfile{frontend/src/app/app.routes.ts}
is what every federated route reaches for:

\begin{lstlisting}[language=javascript]
export function loadRemote<T>(
  remoteName: string, exposedModule: string, exportName: string,
): () => Promise<T> {
  return async () => {
    const loadRemoteModule = (globalThis as any)
      .__SNAPDRAGON_LOAD_REMOTE_MODULE__;
    if (!loadRemoteModule) throw new Error('federation runtime missing');
    const m = await loadRemoteModule({ remoteName, exposedModule });
    return m[exportName];
  };
}
\end{lstlisting}

The exported factory is plugged into \texttt{loadComponent} for every
federated route:

\begin{lstlisting}[language=javascript]
{ path: 'login',       loadComponent: loadRemote('login', './LoginComponent', 'LoginComponent') },
{ path: 'admin',       loadComponent: loadRemote('admin', './AdminShell',     'AdminShell')     },
{ path: 'admin/games', loadComponent: loadRemote('admin', './UploadShell',    'UploadShell')    },
\end{lstlisting}

\section{Per-game federation}

A game remote exposes a \emph{Web Component} (built with
\texttt{createCustomElement} and registered with
\texttt{customElements.define}). The host places it as a normal element:

\begin{lstlisting}[language=html]
@switch (activeRemoteSlug()) {
  @case ('three-wizards') {
    <app-three-wizards-remote
      [levelSlug]="activeLevelSlug()"
      [levelId]="activeLevelId()"
      [launchUid]="activeLaunchUid()" />
  }
  @case ('dodecahedron') {
    <app-dodecahedron-remote [launchUid]="activeLaunchUid()" />
  }
}
\end{lstlisting}

The remote internally builds an iframe URL that points at the
\texttt{wasm-fragment} remote:

\begin{lstlisting}[language=javascript]
this.iframeUrl = `/_remotes/wasm-fragment/${this.levelSlug}/index.html` +
                 `?level_id=${this.levelId}&launch_uid=${this.launchUid}` +
                 `&t=${Date.now()}`;
\end{lstlisting}

The cache-busting \texttt{\&t=...} parameter is intentional — the iframe
HTML itself is small and cheap, but the browser is otherwise prone to
caching it across rebuilds and showing a stale Godot bundle.

\section{Error handling}

When a federated chunk fails to load (host nginx down, remote returned
500, network drop), the host's
\snapfile{AppComponent} subscribes to \texttt{Router.events} and shows a
banner with the human-readable label of the broken route, plus a
\textit{Retry} button. The label lookup uses
\texttt{FEDERATED\_ROUTE\_LABELS} (longest-prefix-first, so
\texttt{/admin/games} resolves before \texttt{/admin}).

\section{Maker remotes (RPG + VN)}

Two newer remotes were added alongside the original four to ship the
\emph{maker} side services as federated Angular projects. They follow
exactly the same pattern as \texttt{login} / \texttt{admin}, but their
backing data-source is a separate Spring Boot side service rather than
the BFF:

\begin{itemize}
  \item \texttt{games-rpg-maker} (kind \texttt{remote-shell}) ---
        mounted at \texttt{/rpg-maker}. Talks to
        \snapfile{rpg-builder/} via the BFF's \texttt{/api/rpg/*} proxy
        path.
  \item \texttt{games-vn-maker} (kind \texttt{remote-shell}) ---
        mounted at \texttt{/vn-maker}. Talks to \snapfile{vn-builder/}
        via \texttt{/api/vn/*}.
\end{itemize}

The federation manifest seed for both lives in
\reposrc{backend/src/main/resources/db/changelog/037-rpg-and-vn-maker-federation.xml}
so the host shell discovers them at bootstrap without a redeploy. The
host nginx config (and the TLS twin --- see chapter~\ref{ch:builddeploy})
needs matching \texttt{location /\_remotes/games-rpg-maker/}
and \texttt{/\_remotes/games-vn-maker/} proxy blocks alongside the
\texttt{/api/rpg/} and \texttt{/api/vn/} sub-service proxies; drift
between the plain and TLS variants produces ``Couldn't reach \emph{X}
--- federated module didn't load.''

\begin{snapcaution}[title={CORP saves WASM iframes}]
WASM games loading inside the federated iframe need
\texttt{Cross-Origin-Resource-Policy: cross-origin} on the
\texttt{wasm-fragment} response \emph{and} the host nginx must strip any
duplicate COOP/COEP headers from the upstream and re-emit a single
canonical set. Missing CORP shows up as
\texttt{net::ERR\_BLOCKED\_BY\_RESPONSE}; duplicates show up as
\texttt{NotSameOriginAfterDefaultedToSameOriginByCoep}. Both are silent
in DevTools' \emph{Console} but show clearly in \emph{Network}'s
\textit{Cross-origin isolation} column.
\end{snapcaution}
